\documentclass{article}
\usepackage[utf8]{inputenc}

\usepackage{amsmath, amssymb, amsthm, amsfonts}
\usepackage{subcaption, enumerate}
\usepackage{bbold}
\usepackage{polynom}
\usepackage{xfrac}
\usepackage{graphicx}
\usepackage{mathtools}
\usepackage{enumitem}

\usepackage{tikz}

\newtheorem{theorem}{Theorem}
\newtheorem{corot}{Corollary}[theorem]
\newtheorem{lemma}{Lemma}
\newtheorem{corol}{Corollary}[lemma]
\newtheorem{proposition}{Proposition}
\newtheorem{corop}{Corollary}[proposition]

\theoremstyle{definition}
\newtheorem{definition}{Definition}
\newtheorem{exampled}{Example}[definition]
\newtheorem{examples}{Example}[section]
\newtheorem*{notation}{Notation}

\newcommand{\dab}
{
	\begin{tikzpicture}[scale=0.04]
		% head
		\draw[thick] (0, 0) circle (3cm);

		% bent arm
		\draw[thick, rotate around={15 : (0, 0)}] (-3, 0) -- (-3.7, -2) -- (4.1, -1.1);

		% extended arm
		\draw[thick, rotate=20] (3, 0) -- (7, 2);
	\end{tikzpicture}
}
\newcommand{\st}{\text{ s.t. }}

\makeatletter
\renewcommand*\env@matrix[1][*\c@MaxMatrixCols c]{%
  \hskip -\arraycolsep
  \let\@ifnextchar\new@ifnextchar
  \array{#1}}
\makeatother

\usepackage{fullpage}

\title{Eigenstuff Day 1 \\ with Dr Nevard}
\author{Laura Gallo}
\date{October 16, 2020}

\begin{document}
\maketitle

\section{Introduction}
For a given $T\in\mathcal{L}(V)$, when can we choose a basis in $V$ such that $T$ has a diagonal matrix? Ie, we want to find some basis such that:
\begin{equation*}
	D=\mathcal{M}(T,\{v_1,v_2,\dots,v_n\})=\left(
		\begin{matrix}
			\lambda_1 & 0 & 0 & 0 \\
			0 & \lambda_2 & 0 & 0 \\
			0 & 0 & \ddots & 0 \\
			0 & 0 & 0 & \lambda_n
		\end{matrix}
	\right)
\end{equation*}
ie, $Tv_k=\lambda_kv_k$ for basis vector $v_k$. This process is called \emph{diagonalisation}, and to do this we must explore eigenvectors and eigenvalues.

\begin{definition}
	Given a linear transform $T:V\rightarrow V$, an \emph{eigenvector} is a nonzero vector $v$ such that $Tv=\lambda v$ for some scalar $\lambda$. $\lambda$ is the corresponding \emph{eigenvalue}.
\end{definition}

\noindent
Note that a linear transformation has an eigenvalue of 0 if and only if it has a nontrivial null space.

\begin{definition}
	Given $T\in\mathcal{L}(V)$, the \emph{eigenspace} corresponding to $\lambda$ is the vector space $\{v:Tv=\lambda v\}$. The proof that this is a vector space is trivial and left as an exercise to the reader. Note that the dimension of an eigenspace can be anywhere between 1 and $\dim V$.
\end{definition}

\section{Examples of Finding Eigenstuff}
\begin{examples}
	Let $V=C^\infty(\mathbb{R})$ (the vector space of all functions that are infinitely differentiable). Let $T(f)=f'$. To find the eigenfunctionss and eigenvalues, we can solve the differental equation:
	$$Tf=f'=\lambda f$$
	The solution to this differential equation is $f=ae^{\lambda x}$, so for a given real eigenvalue $\lambda$, its eigenspace consists of all functions of the form $ae^{\lambda x}$. This means that every real number is an eigenvalue.
\end{examples}

\begin{examples}
	Let $V=\{\text{$C^\infty$ functions of compact support}\}$, $Tf=\int_0^xf(t)dt$. Compact support means that a function $f$ is 0 outside of $[-M, M]$. To find the eigenvectors and eigenvalues, we again solve a differential equation:
	$$Tf=\int_0^xf(t)dt=\lambda f$$
	$$f=\lambda f'$$
	$$f(x)=ae^\frac{x}{\lambda}$$
	However, plugging in 0 into the integral yields 0, which means $a=0$. Since 0 cannot be an eigenvector by definition, this linear transform has no eigenvectors or eigenvalues.
\end{examples}

\begin{examples}
	Let $V=\mathbb{R}^2$, \begin{math}
		T\left(
			\begin{matrix}
				x \\
				y
			\end{matrix}
		\right) = \left(
			\begin{matrix}
				\cos\theta & -\sin\theta \\
				\sin\theta & \cos\theta
			\end{matrix}
		\right) \left(
			\begin{matrix}
				x \\
				y
			\end{matrix}
		\right)
	\end{math}. For $\theta=k\pi$ for some integer $k$, the matrix simplifies to $-I$, so the eigenvalue is -1 and the eigenvectors are all vectors in $\mathbb{R}^2$ except the zero vector. For nontrivial $\theta$, we can do the following to find the eigenvectors and eigenvalues:
	\begin{equation*}
		\left(
			\begin{matrix}
				\cos\theta & -\sin\theta \\
				\sin\theta & \cos\theta
			\end{matrix}
		\right) \left(
			\begin{matrix}
				x \\
				y
			\end{matrix}
		\right)=\lambda \left(
			\begin{matrix}
				x \\
				y
			\end{matrix}
		\right)
	\end{equation*}
	\begin{equation*}
		\left(
			\begin{matrix}
				\cos\theta-\lambda & -\sin\theta \\
				\sin\theta & \cos\theta-\lambda
			\end{matrix}
		\right) \left(
			\begin{matrix}
				x \\
				y
			\end{matrix}
		\right)= \left(
			\begin{matrix}
				0 \\
				0
			\end{matrix}
		\right)
	\end{equation*}
	\begin{align*}
		(\cos\theta-\lambda)x&=\lambda y \\
		(\sin\theta)x&=(\lambda-\cos\theta)y
	\end{align*}
	\begin{align*}
		y&=\left(\frac{\cos\theta-\lambda}{\sin\theta}\right)x \\
		x&=-\left(\frac{\cos\theta-\lambda}{\sin\theta}\right)^2x \\
		0&=\left(1+\left(\frac{\cos\theta-\lambda}{\sin\theta}\right)\right)x
	\end{align*}
	$$x=y=0$$
	Since 0 is not an eigenvector by definition, $T$ has no eigenvectors or eigenvalues.
\end{examples}

\begin{examples}
	Let $V=\mathbb{C}$, $Tv=e^{i\theta}v$. Then every vector is an eigenvector by definition with a corresponding eigenvalue of $e^{i\theta}$.
\end{examples}

\section{Finding Eigenvalues}
We have $Tv=\lambda v$ if and only if $(T-\lambda I)v=0$. This means that the linear transform $T-\lambda I$ has a nontrivial null space, which means that the transform is not invertible and thus its corresponding matrix is singular. Let $A=\mathcal{M}(T)$ for some basis, then $\lambda$ is an eigenvalue if and only if $\det(\lambda I-A)=0$. This equation is the \emph{characteristic equation} of the matrix $A$, and the left hand side is the \emph{characteristic polynomial}.

\begin{examples}
	\begin{math}
		A=\left(
			\begin{matrix}
				2 & 3 \\
				3 & 5
			\end{matrix}
		\right)
	\end{math}
	\begin{equation*}
		\left|
			\begin{matrix}
				\lambda - 2 & -3 \\
				-3 & \lambda - 5
			\end{matrix}
		\right| = (\lambda-2)(\lambda-5)-9=\lambda^2-7\lambda+1
	\end{equation*}
	$$\lambda=\frac{7\pm3\sqrt{5}}{2}$$
\end{examples}

\end{document}

